\section{Distributional Fidelity of the Training Signal} \label{sec:13-introduction-fidelity}

\begin{foldable}
  The restrictions described in the previous section do not merely reduce the quantity of available supervision — they reshape its statistical character.
  A student trained under few-shot constraints, or a surrogate trained on purely synthetic data, or a model fine-tuned on a compressed corpus can only learn from whatever training signal is actually present.
  If that signal fails to faithfully represent the distribution the teacher was trained on, then the student inherits a systematically distorted picture of the teacher's knowledge, regardless of how sophisticated the distillation method itself may be.
  The term \emph{distributional fidelity} names this property: the degree to which the training signal available to the student covers the target distribution in a representative and unbiased way.
  When distributional fidelity is low, distillation fails not because the algorithm is wrong but because the information it needs simply is not present in the signal it receives.
\end{foldable}

\begin{foldable}
  In the model \ac{KD} settings addressed in this thesis — few-shot and data-free — distributional fidelity manifests as a coverage problem on the generative side.
  When a student has no real data and must rely on synthesised images, the generator faces no external pressure to spread its outputs across the full diversity of the input space.
  Instead it gravitates toward a small number of configurations that are easy to produce and reliably elicit confident teacher responses: canonical visual prototypes that sit near the densest, most separable regions of each class's distribution.
  These prototypes are useful, but they represent only a narrow slice of what the teacher knows.
  The teacher's knowledge about less typical instances — unusual viewpoints, ambiguous boundaries between classes, low-contrast or partially occluded examples — remains entirely untapped.
  The student trained on such a corpus achieves good performance on easy, prototypical inputs and generalises poorly elsewhere.
  Distributional fidelity here demands that the synthesised training corpus visit semantically distinct regions across the teacher's full output distribution, not merely the modes the generator finds convenient.
\end{foldable}

\begin{foldable}
  In the dataset distillation setting addressed in this thesis, distributional fidelity takes a different form, with two structurally distinct failure modes.
  The first concerns the scoring of candidate samples.
  Importance-based methods assign each training example a weight that reflects its estimated contribution to a trained model, then select the highest-scoring examples for the distilled corpus.
  At model convergence, however, this scoring procedure is structurally biased: clean samples whose loss has been driven close to zero produce small gradients and receive low importance scores, while noisy, mislabelled, or atypically difficult samples retain large residual gradients and score highly.
  The result is that a naive selection based on these scores fills the distilled corpus disproportionately with hard and often unreliable examples, systematically crowding out the moderate, clean samples that sit near decision boundaries and carry the most robust learning signal for generalization.
  The selected subset is representative of the tail of the training distribution, not the bulk of it.
  The second failure mode arises even if the scoring bias is corrected.
  When samples are selected independently — each judged solely on its own score without regard for what else has already been chosen — redundancy goes unpenalised.
  Two similarly high-scoring samples that represent essentially the same region of the input space will both be selected, consuming budget that could instead cover an entirely unrepresented part of the distribution.
  Correcting the scores does not automatically produce distributional coverage; coverage must be enforced as a separate objective.
\end{foldable}

\begin{foldable}
  Taken together, these failure modes — generative collapse in model \ac{KD} and biased, redundant selection in dataset distillation — are instances of the same underlying problem.
  In every case the training signal that reaches the student does not faithfully represent the target distribution: it over-represents some regions and under-represents or entirely omits others.
  The surface manifestations differ across settings, but the diagnostic question is identical: does the data the student actually trains on give sufficient, representative coverage of the knowledge the teacher possesses?
  Distributional fidelity in the generative sense — diversity of synthetic coverage — and distributional fidelity in the selection sense — unbiased, non-redundant subset coverage — are the same principle operating in different regimes.
  Recognising this unifies what might otherwise appear to be three separate technical problems.
\end{foldable}

\begin{foldable}
  Treating distributional fidelity as a first-class design objective, rather than a hoped-for by-product of running any reasonable method, is the commitment that motivates all three contributions of this thesis.
  Each contribution addresses the failure mode most relevant to its setting, but all three are animated by the same conviction: the quality of distillation is bounded not by the sophistication of the algorithm but by the representativeness of the signal from which it learns.
  \S\ref{sec:14-introduction-contributions} describes how each contribution operationalises this principle in concrete, experimentally validated form.
\end{foldable}
